\documentclass[11pt]{article}
\usepackage[utf8]{inputenc}
\usepackage[spanish, es-tabla]{babel}
\usepackage{amsmath, amssymb}
\usepackage{booktabs}
\usepackage{enumitem}
\usepackage[margin=2.5cm]{geometry}

\setlength{\parindent}{0pt}
\setlength{\parskip}{0.65em}

\title{Econometría ICS-294\\Guía de ejercicios: Clase 9}
\author{Francisco Alfaro Medina}
\date{}

\begin{document}
\maketitle

\section*{Objetivo}

Practicar la interpretación de modelos con logaritmos, variables cuadráticas, interacciones y cambios de escala.

\section*{Ejercicios}

\begin{enumerate}[label=\textbf{Ejercicio \arabic*.}, leftmargin=*]

\item \textbf{Modelo nivel-nivel.}

Considere:
\[
\widehat{salario}=500+80\,educ.
\]

\begin{enumerate}[label=\alph*)]
    \item Interprete el coeficiente de \(educ\).
    \item Prediga el salario para \(educ=12\).
    \item ¿En qué unidades está medida la pendiente?
\end{enumerate}

\item \textbf{Modelo log-nivel.}

Considere:
\[
\widehat{\log(salario)}=1.5+0.07\,educ.
\]

\begin{enumerate}[label=\alph*)]
    \item Interprete el coeficiente de \(educ\).
    \item ¿Cuál es el cambio porcentual aproximado en salario por un año adicional de educación?
    \item ¿Por qué esta interpretación es distinta del modelo nivel-nivel?
\end{enumerate}

\item \textbf{Modelo log-log.}

Considere:
\[
\widehat{\log(consumo)}=0.8+0.65\,\log(ingreso).
\]

\begin{enumerate}[label=\alph*)]
    \item Interprete el coeficiente \(0.65\).
    \item Si el ingreso aumenta en 10\%, ¿cuánto cambia aproximadamente el consumo?
    \item ¿Qué significa decir que el coeficiente es una elasticidad?
\end{enumerate}

\item \textbf{Modelo nivel-log.}

Considere:
\[
\widehat{gasto}=200+15\,\log(ingreso).
\]

\begin{enumerate}[label=\alph*)]
    \item Interprete el coeficiente de \(\log(ingreso)\).
    \item Si el ingreso aumenta en 1\%, ¿cuánto cambia aproximadamente el gasto?
    \item Si el ingreso aumenta en 10\%, ¿cuánto cambia aproximadamente el gasto?
\end{enumerate}

\item \textbf{Término cuadrático.}

Considere:
\[
\widehat{salario}=100+20\,exper-0.4\,exper^2.
\]

\begin{enumerate}[label=\alph*)]
    \item Calcule el efecto parcial de experiencia.
    \item Evalúe el efecto parcial cuando \(exper=5\), \(exper=15\) y \(exper=25\).
    \item ¿A partir de qué nivel de experiencia el efecto se vuelve negativo?
\end{enumerate}

\item \textbf{Interacción.}

Considere:
\[
\widehat{salario}=300+50\,educ+20\,exper+4(educ\cdot exper).
\]

\begin{enumerate}[label=\alph*)]
    \item Calcule el efecto parcial de educación.
    \item Evalúe el efecto de educación cuando \(exper=0\), \(exper=5\) y \(exper=10\).
    \item Explique qué significa que exista una interacción positiva.
\end{enumerate}

\item \textbf{Interacción con variable dummy.}

Considere:
\[
\widehat{salario}=400+60\,educ-100\,mujer+10(educ\cdot mujer).
\]

\begin{enumerate}[label=\alph*)]
    \item ¿Cuál es el retorno a la educación para hombres?
    \item ¿Cuál es el retorno a la educación para mujeres?
    \item Interprete el coeficiente de la interacción.
\end{enumerate}

\item \textbf{Reescalamiento de \(Y\).}

Suponga que:
\[
\widehat{salario\_pesos}=500000+80000\,educ.
\]

Si se mide salario en miles de pesos:
\begin{enumerate}[label=\alph*)]
    \item Escriba la nueva ecuación estimada.
    \item ¿Cambia la interpretación económica?
    \item ¿Cambia el \(R^2\)?
\end{enumerate}

\item \textbf{Reescalamiento de \(X\).}

Suponga que:
\[
\widehat{salario}=300+120\,educ,
\]
donde \(educ\) está medido en años. Si ahora se mide educación en semestres:

\begin{enumerate}[label=\alph*)]
    \item ¿Cuál es la nueva pendiente?
    \item Interprete la nueva pendiente.
    \item ¿Cambia el efecto de un año adicional de educación?
\end{enumerate}

\end{enumerate}

\section*{Preguntas de cierre}

\begin{enumerate}[label=\arabic*.]
    \item ¿Cuándo un coeficiente se interpreta como elasticidad?
    \item ¿Qué permite capturar un término cuadrático?
    \item ¿Qué permite capturar una interacción?
    \item ¿Por qué cambiar unidades no cambia el contenido económico del modelo?
\end{enumerate}

\end{document}
